\section{Outcome (g): Law of total probability and Bayes' theorem}

\subsection{What the outcome asks}

The exam partitions a population into classes (risk classes of drivers, smokers and
non-smokers, suppliers of a part, diseased and healthy patients) and gives the proportion of
the population in each class. These proportions are the \emph{prior} probabilities. Each
class has its own conditional rate of some outcome (an accident, a death, a defect, a positive
test), and the question runs in one of two directions. The forward direction is the law of
total probability: combine the class proportions and the class rates to find the overall rate
of the outcome. The reverse direction is Bayes' theorem: given that the outcome occurred,
find the probability that the individual came from a particular class. Almost every question
on this outcome adds one twist: two years of experience that are independent only within a
class, a class proportion left as an unknown to be solved from a given overall rate, or
conditioning on the outcome \emph{not} occurring.

\subsection{Law of total probability}

Events $B_1, B_2, \dots, B_n$ form a \emph{partition} of the sample space if they are
mutually exclusive ($B_i \cap B_j = \emptyset$ for $i \neq j$) and exhaustive
($B_1 \cup \dots \cup B_n = S$). Every outcome lies in exactly one $B_i$. The simplest
partition is $\{B, \cc{B}\}$.

\begin{keyfact}[Law of total probability]
If $B_1, \dots, B_n$ partition the sample space and $\Prob(B_i) > 0$ for every $i$, then for
any event $A$,
\[
\Prob(A) \;=\; \sum_{i=1}^{n} \Prob(A \mid B_i)\,\Prob(B_i)
\;=\; \Prob(A \mid B_1)\Prob(B_1) + \dots + \Prob(A \mid B_n)\Prob(B_n).
\]
Two-event case:
\[
\Prob(A) \;=\; \Prob(A \mid B)\,\Prob(B) + \Prob(A \mid \cc{B})\,\Prob(\cc{B}).
\]
In words: the overall rate of $A$ is the weighted average of the class rates of $A$, the
weights being the class proportions.
\end{keyfact}

\textbf{Proof.} Because the $B_i$ partition $S$, the event $A$ decomposes into the disjoint
pieces $A \cap B_1, \dots, A \cap B_n$, so $\Prob(A) = \sum_i \Prob(A \cap B_i)$ by
finite additivity. By the multiplication rule, $\Prob(A \cap B_i) = \Prob(A \mid B_i)\Prob(B_i)$.
Substituting gives the formula. \qed

\textbf{Tree view.} The formula is a two-level tree. The first level branches on the class,
with branch probabilities $\Prob(B_i)$; the second level branches on whether $A$ occurs,
with branch probabilities $\Prob(A \mid B_i)$ and $\Prob(\cc{A} \mid B_i) = 1 - \Prob(A \mid B_i)$.
Multiply along each path to get a joint probability, then add the paths that end in $A$.
\begin{itemize}[nosep,leftmargin=2em]
  \item $B_1$ with probability $\Prob(B_1)$
  \begin{itemize}[nosep]
    \item $A$: path probability $\Prob(B_1)\Prob(A \mid B_1)$
    \item $\cc{A}$: path probability $\Prob(B_1)\bigl(1 - \Prob(A \mid B_1)\bigr)$
  \end{itemize}
  \item $B_2$ with probability $\Prob(B_2)$
  \begin{itemize}[nosep]
    \item $A$: path probability $\Prob(B_2)\Prob(A \mid B_2)$
    \item $\cc{A}$: path probability $\Prob(B_2)\bigl(1 - \Prob(A \mid B_2)\bigr)$
  \end{itemize}
  \item $\vdots$
  \item $B_n$ with probability $\Prob(B_n)$
  \begin{itemize}[nosep]
    \item $A$: path probability $\Prob(B_n)\Prob(A \mid B_n)$
    \item $\cc{A}$: path probability $\Prob(B_n)\bigl(1 - \Prob(A \mid B_n)\bigr)$
  \end{itemize}
\end{itemize}
The $2n$ path probabilities sum to 1. $\Prob(A)$ is the sum of the $n$ paths ending in $A$;
$\Prob(\cc{A})$ is the sum of the $n$ paths ending in $\cc{A}$.

\subsection{Bayes' theorem}

Bayes' theorem reverses the direction of conditioning. The data give
$\Prob(A \mid B_i)$ (the rate of the outcome within a class); the question asks for
$\Prob(B_i \mid A)$ (the probability of the class given the outcome).

\begin{keyfact}[Bayes' theorem]
Two events, with $\Prob(A) > 0$ and $0 < \Prob(B) < 1$:
\[
\Prob(B \mid A) \;=\; \frac{\Prob(A \mid B)\,\Prob(B)}{\Prob(A)}
\;=\; \frac{\Prob(A \mid B)\,\Prob(B)}{\Prob(A \mid B)\,\Prob(B) + \Prob(A \mid \cc{B})\,\Prob(\cc{B})}.
\]
Partition $B_1, \dots, B_n$:
\[
\Prob(B_k \mid A) \;=\; \frac{\Prob(A \mid B_k)\,\Prob(B_k)}{\displaystyle\sum_{i=1}^{n} \Prob(A \mid B_i)\,\Prob(B_i)}.
\]
In words:
\[
\text{posterior} \;=\; \frac{\text{prior} \times \text{likelihood}}{\text{sum over all classes of (prior} \times \text{likelihood)}}.
\]
Here the prior is $\Prob(B_k)$, the likelihood is $\Prob(A \mid B_k)$, and the denominator is
$\Prob(A)$ from the law of total probability.
\end{keyfact}

\textbf{Derivation.} By the definition of conditional probability,
$\Prob(B_k \mid A) = \Prob(A \cap B_k)/\Prob(A)$. The numerator is
$\Prob(A \mid B_k)\Prob(B_k)$ by the multiplication rule, and the denominator is
$\sum_i \Prob(A \mid B_i)\Prob(B_i)$ by the law of total probability. \qed

Two consequences are worth keeping in view. First, the posterior probabilities
$\Prob(B_i \mid A)$ for $i = 1, \dots, n$ sum to 1, because the numerators sum to the
denominator. Second, the ratio of two posteriors is
\[
\frac{\Prob(B_j \mid A)}{\Prob(B_k \mid A)} = \frac{\Prob(B_j)\,\Prob(A \mid B_j)}{\Prob(B_k)\,\Prob(A \mid B_k)},
\]
because the denominator cancels. The observation $A$ multiplies each class's prior odds by
that class's likelihood.

\subsection{The table method}

Every question on this outcome can be organised in one table with one row per class. The
steps:
\begin{enumerate}[nosep,leftmargin=2em]
  \item List the classes $B_i$ (they must cover everyone and not overlap).
  \item Write the prior $\Prob(B_i)$ in column 2. The column must sum to 1.
  \item Write the likelihood $\Prob(A \mid B_i)$ in column 3, where $A$ is the event
        actually observed. If the observed event is ``no accident'', column 3 holds
        $1 - (\text{accident rate})$.
  \item Multiply across: joint $\Prob(A \cap B_i) = \Prob(B_i)\,\Prob(A \mid B_i)$ in column 4.
  \item Sum column 4. The sum is $\Prob(A)$ (law of total probability).
  \item Divide each joint by the column sum: posterior $\Prob(B_i \mid A)$ in column 5
        (Bayes' theorem). Column 5 must sum to 1.
\end{enumerate}

\textbf{Worked table.} A health insurer classifies its policyholders as smokers (20\%),
former smokers (30\%), and never-smokers (50\%). The probability of a claim in a year is 0.30
for a smoker, 0.20 for a former smoker, and 0.10 for a never-smoker. A randomly chosen
policyholder files a claim. Find the probability that the policyholder is a smoker.

\begin{center}
\begin{tabular}{@{}lcccc@{}}
\toprule
Class $B_i$ & Prior $\Prob(B_i)$ & Likelihood $\Prob(A \mid B_i)$ & Joint $\Prob(A \cap B_i)$ & Posterior $\Prob(B_i \mid A)$ \\
\midrule
Smoker        & 0.20 & 0.30 & $0.20 \times 0.30 = 0.06$ & $0.06/0.17 = 0.3529$ \\
Former smoker & 0.30 & 0.20 & $0.30 \times 0.20 = 0.06$ & $0.06/0.17 = 0.3529$ \\
Never-smoker  & 0.50 & 0.10 & $0.50 \times 0.10 = 0.05$ & $0.05/0.17 = 0.2941$ \\
\midrule
Total         & 1.00 &      & $\Prob(A) = 0.17$          & 1.0000 \\
\bottomrule
\end{tabular}
\end{center}

The overall claim rate is $\Prob(A) = 0.17$, and $\Prob(\text{smoker} \mid \text{claim}) = 0.06/0.17 = 0.3529$.
Note that the prior for smokers was 0.20; observing a claim raised it to 0.35.

\textbf{Natural frequencies.} When the priors are proportions of a population, the same
computation can be done by counting in a hypothetical population. Take 1000 policyholders.
\begin{itemize}[nosep,leftmargin=2em]
  \item Smokers: $0.20 \times 1000 = 200$ people, of whom $0.30 \times 200 = 60$ file a claim.
  \item Former smokers: $0.30 \times 1000 = 300$ people, of whom $0.20 \times 300 = 60$ file a claim.
  \item Never-smokers: $0.50 \times 1000 = 500$ people, of whom $0.10 \times 500 = 50$ file a claim.
\end{itemize}
Total claimants: $60 + 60 + 50 = 170$. Of the 170 claimants, 60 are smokers, so
$\Prob(\text{smoker} \mid \text{claim}) = 60/170 = 0.3529$. The counts are the joint column
multiplied by 1000; the arithmetic is the same, but the counts are harder to get wrong under
time pressure. Use 10\,000 or 100\,000 when the rates are small (a prevalence of 0.001, say)
so that every count is a whole number.

\subsection{Worked examples}

\begin{example}[Three risk classes of drivers]
An auto insurer classifies its drivers as standard, preferred, or ultra-preferred. Of the
insurer's drivers, 60\% are standard, 30\% are preferred, and 10\% are ultra-preferred. In a
given year, the probability of an accident is 0.10 for a standard driver, 0.05 for a
preferred driver, and 0.01 for an ultra-preferred driver. A driver chosen at random has an
accident during the year. Calculate the probability that the driver is ultra-preferred.
\par\smallskip\noindent (A) 0.001\qquad (B) 0.010\qquad (C) 0.013\qquad (D) 0.076\qquad (E) 0.132
\end{example}
\begin{solution}
Let $A$ be the event of an accident. Build the table.
\begin{center}
\begin{tabular}{@{}lccc@{}}
\toprule
Class & Prior & Likelihood & Joint \\
\midrule
Standard        & 0.60 & 0.10 & 0.060 \\
Preferred       & 0.30 & 0.05 & 0.015 \\
Ultra-preferred & 0.10 & 0.01 & 0.001 \\
\midrule
Total           & 1.00 &      & $\Prob(A) = 0.076$ \\
\bottomrule
\end{tabular}
\end{center}
\[
\Prob(\text{ultra} \mid A) = \frac{0.001}{0.076} = 0.0132.
\]
Answer (C). Choice (A) is the joint probability, not the posterior; choice (D) is $\Prob(A)$.
Note the base-rate effect: ultra-preferred drivers are 10\% of the population but only 1.3\%
of the drivers who have accidents.
\end{solution}

\begin{example}[Medical test with low prevalence]
A blood test for a disease has sensitivity 0.95 (the probability of a positive result given
that the disease is present) and specificity 0.90 (the probability of a negative result given
that the disease is absent). The disease is present in 1\% of the population. A person chosen
at random tests positive. Calculate the probability that the person has the disease.
\par\smallskip\noindent (A) 0.009\qquad (B) 0.088\qquad (C) 0.095\qquad (D) 0.475\qquad (E) 0.950
\end{example}
\begin{solution}
Let $D$ be disease present and $T$ a positive test. The likelihood of a positive test in the
healthy class is the false-positive rate $1 - \text{specificity} = 0.10$.
\begin{center}
\begin{tabular}{@{}lccc@{}}
\toprule
Class & Prior & $\Prob(T \mid \cdot)$ & Joint \\
\midrule
Disease    & 0.01 & 0.95 & 0.0095 \\
No disease & 0.99 & 0.10 & 0.0990 \\
\midrule
Total      & 1.00 &      & $\Prob(T) = 0.1085$ \\
\bottomrule
\end{tabular}
\end{center}
\[
\Prob(D \mid T) = \frac{0.0095}{0.1085} = 0.0876.
\]
Answer (B). In natural frequencies with 10\,000 people: 100 have the disease and 95 of them
test positive; 9\,900 do not and 990 of them test positive; $95/(95 + 990) = 0.0876$. A test
that is 95\% sensitive and 90\% specific still leaves a positive result with under a 9\%
chance of disease, because the healthy class is 99 times larger than the diseased class and
its false positives dominate. Choice (E) confuses $\Prob(T \mid D)$ with $\Prob(D \mid T)$.
\end{solution}

\begin{example}[Claim in year two given no claim in year one]
An insurer has two classes of policyholders. Low-risk policyholders make up 70\% of the
insured population and have probability 0.10 of a claim in any given year. High-risk
policyholders make up the remaining 30\% and have probability 0.40 of a claim in any given
year. For each policyholder, claims in different years are independent given the
policyholder's class. A policyholder chosen at random has no claim in the first year.
Calculate the probability that the policyholder has a claim in the second year.
\par\smallskip\noindent (A) 0.135\qquad (B) 0.167\qquad (C) 0.190\qquad (D) 0.210\qquad (E) 0.250
\end{example}
\begin{solution}
Let $N_1$ be no claim in year 1 and $C_2$ a claim in year 2. Then
$\Prob(C_2 \mid N_1) = \Prob(N_1 \cap C_2)/\Prob(N_1)$, and both numerator and denominator
are computed by total probability over the classes. Within a class the years are independent,
so $\Prob(N_1 \cap C_2 \mid \text{class}) = \Prob(N_1 \mid \text{class})\,\Prob(C_2 \mid \text{class})$.
\begin{center}
\begin{tabular}{@{}lcccc@{}}
\toprule
Class & Prior & $\Prob(N_1 \mid \cdot)$ & Joint $\Prob(N_1 \cap \cdot)$ & $\Prob(N_1 \cap C_2 \cap \cdot)$ \\
\midrule
Low  & 0.70 & 0.90 & 0.630 & $0.630 \times 0.10 = 0.063$ \\
High & 0.30 & 0.60 & 0.180 & $0.180 \times 0.40 = 0.072$ \\
\midrule
Total & 1.00 & & $\Prob(N_1) = 0.810$ & $\Prob(N_1 \cap C_2) = 0.135$ \\
\bottomrule
\end{tabular}
\end{center}
\[
\Prob(C_2 \mid N_1) = \frac{0.135}{0.810} = 0.1667.
\]
Answer (B). An equivalent route is to update the class probabilities first:
$\Prob(\text{low} \mid N_1) = 0.630/0.810 = 7/9$ and $\Prob(\text{high} \mid N_1) = 2/9$; then
$\Prob(C_2 \mid N_1) = (7/9)(0.10) + (2/9)(0.40) = 0.1667$. Choice (C) is the unconditional
one-year claim rate $0.7(0.10) + 0.3(0.40) = 0.19$; the clean first year shifts weight toward
the low-risk class and lowers the second-year rate. The years are \emph{not} unconditionally
independent, so $\Prob(C_2 \mid N_1) \neq \Prob(C_2)$.
\end{solution}

\begin{example}[Unknown class proportion]
An insurer issues policies to two types of policyholders. Type A policyholders have
probability 0.05 of a claim in a year; type B policyholders have probability 0.20. The
insurer's overall claim rate is 0.11 per policyholder per year. A policyholder chosen at
random files a claim. Calculate the probability that the policyholder is type B.
\par\smallskip\noindent (A) 0.40\qquad (B) 0.60\qquad (C) 0.67\qquad (D) 0.73\qquad (E) 0.80
\end{example}
\begin{solution}
Let $x = \Prob(\text{A})$, so $\Prob(\text{B}) = 1 - x$. The law of total probability gives
\[
0.11 = 0.05x + 0.20(1 - x) = 0.20 - 0.15x
\quad\Longrightarrow\quad
x = \frac{0.09}{0.15} = 0.60.
\]
So 60\% are type A and 40\% are type B. Now apply Bayes' theorem; the denominator is the
given overall rate 0.11, so there is nothing further to add up:
\[
\Prob(\text{B} \mid \text{claim}) = \frac{0.40 \times 0.20}{0.11} = \frac{0.08}{0.11} = 0.7273.
\]
Answer (D). Choice (A) is the prior for type B; choice (B) is the proportion of type A.
\end{solution}

\begin{example}[Classes given as a ratio]
A company has twice as many male employees as female employees. In a given year, 6\% of male
employees and 9\% of female employees file a disability claim. An employee chosen at random
files a disability claim. Calculate the probability that the employee is female.
\par\smallskip\noindent (A) 0.333\qquad (B) 0.400\qquad (C) 0.429\qquad (D) 0.500\qquad (E) 0.600
\end{example}
\begin{solution}
``Twice as many male as female'' means the ratio male : female is $2 : 1$, so the priors are
$\Prob(M) = 2/3$ and $\Prob(F) = 1/3$. (The same information stated as odds ``$2$ to $1$
in favour of male'' converts the same way: odds $a : b$ means probabilities $a/(a+b)$ and
$b/(a+b)$.)
\begin{center}
\begin{tabular}{@{}lccc@{}}
\toprule
Class & Prior & Likelihood & Joint \\
\midrule
Male   & $2/3$ & 0.06 & 0.04 \\
Female & $1/3$ & 0.09 & 0.03 \\
\midrule
Total  & 1 & & 0.07 \\
\bottomrule
\end{tabular}
\end{center}
\[
\Prob(F \mid \text{claim}) = \frac{0.03}{0.07} = 0.4286.
\]
Answer (C). Using natural frequencies with 300 employees: 200 males with 12 claims, 100
females with 9 claims; $9/21 = 0.4286$. Choice (A) is the prior. Choice (E) comes from
using the wrong ratio (treating female as the larger class).
\end{solution}

\begin{example}[Which supplier]
A manufacturer buys a component from three suppliers. Supplier A provides 50\% of the
components, supplier B provides 30\%, and supplier C provides 20\%. The proportion of
defective components is 2\% for supplier A, 3\% for supplier B, and 5\% for supplier C. A
component chosen at random from inventory is found to be defective. Calculate the probability
that it came from supplier C.
\par\smallskip\noindent (A) 0.200\qquad (B) 0.310\qquad (C) 0.345\qquad (D) 0.400\qquad (E) 0.500
\end{example}
\begin{solution}
\begin{center}
\begin{tabular}{@{}lcccc@{}}
\toprule
Supplier & Prior & Likelihood & Joint & Posterior \\
\midrule
A & 0.50 & 0.02 & 0.010 & $0.010/0.029 = 0.345$ \\
B & 0.30 & 0.03 & 0.009 & $0.009/0.029 = 0.310$ \\
C & 0.20 & 0.05 & 0.010 & $0.010/0.029 = 0.345$ \\
\midrule
Total & 1.00 & & 0.029 & 1.000 \\
\bottomrule
\end{tabular}
\end{center}
$\Prob(\text{C} \mid \text{defective}) = 0.010/0.029 = 0.3448$. Answer (C). Choice (B) is the
posterior for supplier B, a distractor for a candidate who reads the wrong row. In natural
frequencies with 1000 components: A supplies 500 with 10 defects, B supplies 300 with 9,
C supplies 200 with 10; $10/29$.
\end{solution}

\begin{example}[Conditioning on the complement]
An insurer's drivers fall into three age groups: 30\% are young, 50\% are middle-aged, and
20\% are senior. The probability of an accident in a year is 0.20 for a young driver, 0.08
for a middle-aged driver, and 0.12 for a senior driver. A driver chosen at random has no
accident during the year. Calculate the probability that the driver is young.
\par\smallskip\noindent (A) 0.240\qquad (B) 0.274\qquad (C) 0.300\qquad (D) 0.484\qquad (E) 0.516
\end{example}
\begin{solution}
The observed event is $\cc{A}$, no accident. The likelihood in each row is
$\Prob(\cc{A} \mid \text{class}) = 1 - \Prob(A \mid \text{class})$, in every row.
\begin{center}
\begin{tabular}{@{}lccc@{}}
\toprule
Class & Prior & $\Prob(\cc{A} \mid \cdot)$ & Joint \\
\midrule
Young  & 0.30 & 0.80 & 0.240 \\
Middle & 0.50 & 0.92 & 0.460 \\
Senior & 0.20 & 0.88 & 0.176 \\
\midrule
Total  & 1.00 & & $\Prob(\cc{A}) = 0.876$ \\
\bottomrule
\end{tabular}
\end{center}
\[
\Prob(\text{young} \mid \cc{A}) = \frac{0.240}{0.876} = 0.2740.
\]
Answer (B). Choice (D), 0.484, is $\Prob(\text{young} \mid A) = 0.060/0.124$, the answer to
the opposite question; choice (E) is $1 - 0.484$, which is $\Prob(\text{not young} \mid A)$
and not $\Prob(\text{young} \mid \cc{A})$. A clean year lowers the probability of the young
class from 0.30 to 0.27, a smaller shift than an accident produces, because a clean year is
common in every class.
\end{solution}

\begin{example}[Observation spanning two years]
An insurer classifies 80\% of its policyholders as good risks and 20\% as bad risks. In each
year a good risk has probability 0.05 of at least one claim and a bad risk has probability
0.25. Given the class, claim activity in different years is independent. A policyholder
chosen at random has at least one claim during a two-year period. Calculate the probability
that the policyholder is a bad risk.
\par\smallskip\noindent (A) 0.200\qquad (B) 0.400\qquad (C) 0.529\qquad (D) 0.556\qquad (E) 0.714
\end{example}
\begin{solution}
Let $E$ be the event of at least one claim in two years. Within a class,
$\Prob(E \mid \text{class}) = 1 - (1 - p)^2$ where $p$ is the class's one-year rate, using
conditional independence.
\begin{center}
\begin{tabular}{@{}lccc@{}}
\toprule
Class & Prior & $\Prob(E \mid \cdot)$ & Joint \\
\midrule
Good & 0.80 & $1 - 0.95^2 = 0.0975$ & 0.0780 \\
Bad  & 0.20 & $1 - 0.75^2 = 0.4375$ & 0.0875 \\
\midrule
Total & 1.00 & & $\Prob(E) = 0.1655$ \\
\bottomrule
\end{tabular}
\end{center}
\[
\Prob(\text{bad} \mid E) = \frac{0.0875}{0.1655} = 0.5287.
\]
Answer (C). Choice (D), 0.556, is $\Prob(\text{bad} \mid \text{claim in a single year}) = 0.05/0.09$;
the likelihood must match the event actually observed, which here covers two years.
\end{solution}

\subsection{Traps}

\begin{trap}[Answering $\Prob(A \mid B)$ when $\Prob(B \mid A)$ is asked]
The data always give the rate of the outcome within each class, $\Prob(A \mid B_i)$. The
question almost always asks for the class given the outcome, $\Prob(B_i \mid A)$. These are
different numbers: in the medical-test example, $\Prob(T \mid D) = 0.95$ while
$\Prob(D \mid T) = 0.088$. Before computing, write down in symbols which event is observed
(goes to the right of the bar) and which is asked for (goes to the left). The answer choices
routinely include the likelihood as a distractor.
\end{trap}

\begin{trap}[Forgetting a class in the denominator]
The denominator of Bayes' theorem is $\Prob(A)$, the sum of prior times likelihood over
\emph{every} class. Omitting a class (typically the third class in a three-class problem, or
the ``no disease'' class in a test problem) makes the denominator too small and the posterior
too large. Check that the prior column sums to 1 before summing the joint column. If the
priors given in the problem do not sum to 1, the remaining mass belongs to a class the
problem describes in words (``the remaining policyholders'') and that class must be in the
table.
\end{trap}

\begin{trap}[Using the overall rate as a likelihood]
The likelihood column holds class-specific rates $\Prob(A \mid B_i)$. The overall rate
$\Prob(A)$ belongs in the total row, not in any class row. When a problem gives the overall
rate together with all but one of the class rates or priors, the overall rate is there so
that the missing quantity can be solved from the law of total probability, as in the
unknown-proportion example. Once solved, the given overall rate is the Bayes denominator; do
not recompute it from the table and get a different number through rounding.
\end{trap}

\begin{trap}[Conditional independence is given the class, not unconditional]
In multi-year problems, ``claims in different years are independent'' means independent
\emph{within each class}: $\Prob(C_1 \cap C_2 \mid B_i) = \Prob(C_1 \mid B_i)\Prob(C_2 \mid B_i)$.
It does not mean $\Prob(C_1 \cap C_2) = \Prob(C_1)\Prob(C_2)$. Because the population mixes
classes with different rates, a claim in year one raises the probability of the high-risk
class and therefore raises the probability of a claim in year two:
$\Prob(C_2 \mid C_1) > \Prob(C_2)$. A candidate who treats the years as unconditionally
independent answers with the marginal one-year rate, which is always among the choices.
Multiply the two years' probabilities inside each row of the table, never across the total
row.
\end{trap}

\begin{trap}[Conditioning on $\cc{A}$ requires $1 - \text{likelihood}$ in every row]
When the observed event is ``no accident'', ``tested negative'', or ``survived'', the
likelihood in every row is $1 - \Prob(A \mid B_i)$, and the denominator is
$\Prob(\cc{A}) = 1 - \Prob(A)$. Two wrong shortcuts appear in the choices:
using the accident rates unchanged (which gives $\Prob(B_i \mid A)$), and computing
$1 - \Prob(B_i \mid A)$ (which is $\Prob(\cc{B_i} \mid A)$, the complement on the wrong
side of the bar). Neither equals $\Prob(B_i \mid \cc{A})$. For a negative medical test the
likelihoods are $1 - \text{sensitivity}$ in the disease row and the specificity in the
healthy row.
\end{trap}

\begin{trap}[Normalising by the wrong total]
The posterior is the joint probability divided by the sum of the joint column, not by the
prior, not by the likelihood, and not by 1. Dividing the joint by the class prior returns the
likelihood; dividing by the likelihood returns the prior; leaving it undivided gives
$\Prob(A \cap B_i)$, which the choices usually include (0.001 in the driver example, 0.018 in
a test example). The posterior column must sum to 1; if it does not, the divisor is wrong.
\end{trap}

\subsection{Practice problems}

\begin{practice}
An insurer classifies its policyholders as low risk (45\%), medium risk (35\%), and high
risk (20\%). The probability of a claim in a year is 0.02 for low risk, 0.05 for medium
risk, and 0.10 for high risk. A policyholder chosen at random files a claim. Calculate the
probability that the policyholder is high risk.
\par\smallskip\noindent (A) 0.200\qquad (B) 0.387\qquad (C) 0.430\qquad (D) 0.465\qquad (E) 0.500
\end{practice}

\begin{practice}
A screening test for a condition has sensitivity 0.90 and specificity 0.95. The condition
is present in 2\% of the population screened. A person screened at random tests positive.
Calculate the probability that the person has the condition.
\par\smallskip\noindent (A) 0.018\qquad (B) 0.067\qquad (C) 0.269\qquad (D) 0.474\qquad (E) 0.900
\end{practice}

\begin{practice}
An insurer's policyholders are 60\% class A, with annual claim probability 0.10, and 40\%
class B, with annual claim probability 0.30. Given the class, claims in different years are
independent. A policyholder chosen at random has a claim in year one. Calculate the
probability that the policyholder has a claim in year two.
\par\smallskip\noindent (A) 0.180\qquad (B) 0.200\qquad (C) 0.233\qquad (D) 0.250\qquad (E) 0.300
\end{practice}

\begin{practice}
A factory produces a part on two machines. Parts from machine 1 are defective with
probability 0.02 and parts from machine 2 are defective with probability 0.05. The
proportion of all parts that are defective is 0.032. A part chosen at random is defective.
Calculate the probability that it was produced on machine 2.
\par\smallskip\noindent (A) 0.400\qquad (B) 0.500\qquad (C) 0.600\qquad (D) 0.625\qquad (E) 0.750
\end{practice}

\begin{practice}
An auto insurer has three times as many preferred drivers as standard drivers. In a year,
a standard driver has an accident with probability 0.12 and a preferred driver with
probability 0.04. A driver chosen at random has an accident. Calculate the probability that
the driver is standard.
\par\smallskip\noindent (A) 0.250\qquad (B) 0.333\qquad (C) 0.500\qquad (D) 0.667\qquad (E) 0.750
\end{practice}

\begin{practice}
Of a life insurer's policyholders, 25\% are smokers and 75\% are non-smokers. The
probability of dying within the year is 0.05 for a smoker and 0.01 for a non-smoker. A
policyholder chosen at random survives the year. Calculate the probability that the
policyholder is a smoker.
\par\smallskip\noindent (A) 0.238\qquad (B) 0.242\qquad (C) 0.250\qquad (D) 0.375\qquad (E) 0.625
\end{practice}

\begin{practice}
A distributor stocks a component from four suppliers. Supplier A provides 40\% of the
stock with a 1\% defect rate, supplier B provides 30\% with a 2\% defect rate, supplier C
provides 20\% with a 3\% defect rate, and supplier D provides 10\% with a 5\% defect rate.
A component chosen at random is defective. Calculate the probability that it came from
supplier A.
\par\smallskip\noindent (A) 0.100\qquad (B) 0.190\qquad (C) 0.286\qquad (D) 0.400\qquad (E) 0.476
\end{practice}

\begin{practice}
An insurer's policyholders are 30\% young and 70\% not young. The probability that a young
policyholder files a claim in a year is 0.15. The overall probability that a policyholder
files a claim in a year is 0.09. A policyholder chosen at random files no claim during the
year. Calculate the probability that the policyholder is young.
\par\smallskip\noindent (A) 0.064\qquad (B) 0.255\qquad (C) 0.280\qquad (D) 0.300\qquad (E) 0.500
\end{practice}

\subsection{Answers to practice problems}

\begin{enumerate}[nosep,leftmargin=2em]
  \item \textbf{(C)}. Joint probabilities $0.45(0.02) = 0.0090$, $0.35(0.05) = 0.0175$,
        $0.20(0.10) = 0.0200$; total $0.0465$. $\Prob(\text{high} \mid \text{claim}) = 0.0200/0.0465 = 0.430$.
  \item \textbf{(C)}. Joint: $0.02(0.90) = 0.018$ and $0.98(0.05) = 0.049$; total $0.067$.
        $\Prob(D \mid T) = 0.018/0.067 = 0.269$. Choice (A) is the joint, (B) is $\Prob(T)$.
  \item \textbf{(C)}. $\Prob(C_1) = 0.6(0.10) + 0.4(0.30) = 0.18$;
        $\Prob(C_1 \cap C_2) = 0.6(0.10)^2 + 0.4(0.30)^2 = 0.006 + 0.036 = 0.042$;
        $\Prob(C_2 \mid C_1) = 0.042/0.18 = 0.233$. Choice (A) is the unconditional rate.
  \item \textbf{(D)}. With $x = \Prob(\text{machine 1})$: $0.02x + 0.05(1 - x) = 0.032$ gives
        $x = 0.60$. $\Prob(\text{machine 2} \mid \text{def}) = 0.40(0.05)/0.032 = 0.020/0.032 = 0.625$.
  \item \textbf{(C)}. Priors $\Prob(\text{std}) = 1/4$, $\Prob(\text{pref}) = 3/4$. Joint:
        $0.25(0.12) = 0.03$ and $0.75(0.04) = 0.03$; total $0.06$. $\Prob(\text{std} \mid A) = 0.03/0.06 = 0.500$.
  \item \textbf{(B)}. Survival likelihoods $0.95$ and $0.99$. Joint: $0.25(0.95) = 0.2375$ and
        $0.75(0.99) = 0.7425$; total $0.98$. $\Prob(\text{smoker} \mid \text{survive}) = 0.2375/0.98 = 0.242$.
        Choice (E) is $\Prob(\text{smoker} \mid \text{die}) = 0.0125/0.02$; choice (D) is its complement.
  \item \textbf{(B)}. Joint: $0.004$, $0.006$, $0.006$, $0.005$; total $0.021$.
        $\Prob(\text{A} \mid \text{def}) = 0.004/0.021 = 0.190$. Choice (C) is the posterior for B or C.
  \item \textbf{(C)}. First solve for the not-young rate $y$: $0.3(0.15) + 0.7y = 0.09$ gives
        $y = 0.045/0.7 = 0.0643$. Then $\Prob(\text{no claim}) = 1 - 0.09 = 0.91$ and
        $\Prob(\text{young} \mid \text{no claim}) = 0.3(0.85)/0.91 = 0.255/0.91 = 0.280$.
        Choice (A) is $y$; choice (B) is the joint probability.
\end{enumerate}

\subsection{One-page summary}

\begin{itemize}[nosep,leftmargin=2em]
  \item A partition $B_1, \dots, B_n$: mutually exclusive, exhaustive, priors sum to 1.
        The two-class case is $\{B, \cc{B}\}$.
  \item Law of total probability: $\Prob(A) = \sum_i \Prob(A \mid B_i)\Prob(B_i)$.
        Two classes: $\Prob(A) = \Prob(A \mid B)\Prob(B) + \Prob(A \mid \cc{B})\Prob(\cc{B})$.
        The overall rate is the prior-weighted average of the class rates.
  \item Bayes' theorem: $\Prob(B_k \mid A) = \dfrac{\Prob(A \mid B_k)\Prob(B_k)}{\sum_i \Prob(A \mid B_i)\Prob(B_i)}$.
        Posterior $=$ prior $\times$ likelihood, divided by the sum of prior $\times$ likelihood
        over all classes. Posteriors sum to 1.
  \item Table method: (1) one row per class; (2) prior column, sums to 1; (3) likelihood
        column $\Prob(\text{observed event} \mid \text{class})$; (4) joint $=$ prior $\times$
        likelihood; (5) sum the joint column to get $\Prob(\text{observed event})$;
        (6) posterior $=$ joint $/$ column sum.
  \item Natural frequencies: take 1000 or 10\,000 people, count each class, count the
        outcome within each class, and divide the count in the target class by the total
        count with the outcome.
  \item Observed event on the right of the bar, asked-for event on the left. The data give
        $\Prob(A \mid B_i)$; the question asks $\Prob(B_i \mid A)$.
  \item Conditioning on $\cc{A}$ (no claim, negative test, survived): likelihood
        $1 - \Prob(A \mid B_i)$ in every row; denominator $1 - \Prob(A)$.
        Negative test: disease row $1 - \text{sensitivity}$, healthy row $=$ specificity.
  \item Sensitivity $= \Prob(T \mid D)$; specificity $= \Prob(\cc{T} \mid \cc{D})$;
        false-positive rate $= 1 - \text{specificity}$. Low prevalence makes $\Prob(D \mid T)$
        small even for an accurate test.
  \item Two years: independence holds within a class only. Multiply year probabilities inside
        each row: $\Prob(C_1 \cap C_2 \mid B_i) = p_i^2$, $\Prob(N_1 \cap C_2 \mid B_i) = (1 - p_i)p_i$,
        $\Prob(\text{at least one claim} \mid B_i) = 1 - (1 - p_i)^2$. Then
        $\Prob(C_2 \mid C_1) = \Prob(C_1 \cap C_2)/\Prob(C_1)$, each side by total probability.
        Equivalent: update the class posteriors after year one, then average the year-two rates.
  \item Unknown proportion: set $\Prob(B_1) = x$, write the given overall rate as
        $\sum_i \Prob(A \mid B_i)\Prob(B_i)$, solve for $x$, then use the given overall rate as
        the Bayes denominator.
  \item Ratios and odds: ``$k$ times as many in class 1 as class 2'' or odds $k : 1$ means
        priors $k/(k+1)$ and $1/(k+1)$.
  \item Distractors to recognise: the likelihood, the joint probability, the overall rate,
        the prior, the posterior of a neighbouring class, and the complement on the wrong
        side of the bar.
\end{itemize}
